% Generated draft: requires mathematical and source review.
\documentclass[11pt]{article}
\usepackage[margin=25mm]{geometry}
\usepackage{fontspec}
\setmainfont{Latin Modern Roman}
\usepackage{amsmath,amssymb,mathtools,unicode-math}
\setmathfont{Latin Modern Math}
\usepackage{enumitem,needspace,xurl,hyperref}
\setlength{\parindent}{0pt}
\setlength{\parskip}{5pt}
\setlength{\emergencystretch}{3em}
\newcommand{\N}{\mathbb{N}}
\newcommand{\Z}{\mathbb{Z}}
\newcommand{\Q}{\mathbb{Q}}
\newcommand{\R}{\mathbb{R}}
\newcommand{\C}{\mathbb{C}}
\newcommand{\F}{\mathbb{F}}
\newcommand{\A}{\mathbb{A}}
\newcommand{\PP}{\mathbb{P}}
\newcommand{\E}{\mathbb{E}}
\newcommand{\eps}{\varepsilon}
\newcommand{\dd}{\,\mathrm{d}}
\begin{document}
\section*{UnsolvedMath ID: 700819}
\textbf{Source problem code:} catalogue entry
\subsection{Definitions and mathematical statement}
Let \(G\) be a group with identity element \(1_G\). The exponent-\(12\) identity is
\[
\forall g\in G,\qquad g^{12}=1_G.
\]
This identity says that the exponent of \(G\) divides \(12\). A group is locally finite if every finitely generated subgroup is finite; equivalently, for every finite subset \(S\subseteq G\), the subgroup \(\langle S\rangle\) generated by \(S\) is finite.

The question attributed to V. P. Shunkov is whether, for every group \(G\),
\[
\bigl(\forall g\in G,\ g^{12}=1_G\bigr)
\quad\Longrightarrow\quad
\bigl(\forall S\subseteq G,\ |S|<\infty\Longrightarrow |\langle S\rangle|<\infty\bigr).
\]
The source uses the phrase ``exponent \(12\)'' without distinguishing exponent dividing \(12\) from least positive exponent exactly \(12\). The corresponding universal questions are equivalent: any group satisfying the displayed identity embeds in its direct product with a cyclic group of order \(12\), which has exponent exactly \(12\). No finiteness assumption is imposed on \(G\).
\subsection{Short English statement}
Must every finitely generated subgroup of every group satisfying \(g^{12}=1\) for all its elements be finite? This is the question proposed by V. P. Shunkov as Kourovka Notebook Problem 11.127.
\subsection{Sources}
\begin{itemize}
\item Local UnsolvedMath registry, permanent ID \(700819\), source problem code \texttt{catalogue entry}, title ``Kourovka Notebook Problem 11.127''; supplied upstream revision \texttt{d371ee917ac9647fbb06cd6fc529deec6f3dda82}.
\item The supplied record attributes the problem to V. P. Shunkov, proposed in \(1990\), and identifies its source as \emph{The Kourovka Notebook (Unsolved Problems in Group Theory)}, 11th issue (1990), Problem 11.127, with reference to the 21st issue edition (2026), \texttt{arXiv:1401.0300v48}.
\item The supplied collection description names E. I. Khukhro and V. D. Mazurov as editors. The registry labels this problem ``open''; this is supplied source metadata, not an independently verified assessment.
\end{itemize}
\end{document}
